\paragraph{\trRiemann\trSurface。}\label{sub:riemann-surfaces}
至此，我们固定一个\trCorneredBoundary{}的\trExactSymplectic\trManifold{}$M$，它将是此后讨论全程的%
“目标空间”。记 $\varpi=\varpi_M$、$\theta=\theta_M$、$I=I_M$ 为 $M$ 上给定的诸结构。
设 $\widehat{S}$ 是一个带边的\trCompact\trRiemann\trSurface，$\Sigma$ 是边界点的一个有限集，分成
两部分 $\Sigma=\Sigma^-\sqcup\Sigma^+$。我们称 $S=\widehat{S}\setminus\Sigma$ 为一个%
\textbf{\trPointedBoundary\trRiemann\trSurface}（pointed-boundary Riemann surface），而 $\Sigma^-$、
$\Sigma^+$ 的点（它们作为数据的组成部分被保留下来）分别是其无穷远处的\trIncomingPoint{}与%
\trOutgoingPoint。注意 $\widehat{S}$ 可以由 $S$ \trCanonical{}地复原。对最简单的几类这样的%
\trSurface，我们将使用如下记号：
\begin{itemize}
\item $D$ 是 $\CC$ 中的闭单位\trDisk；
\item $H$ 是闭\trUpperHalfPlane，带一个无穷远处的\trIncomingPoint；
\item $\overline{H}$ 与 $H$ 是同一个\trSurface，只是其无穷远处的那个点被当作\trOutgoingPoint；
\item $Z=\RR\times[0;1]$，带坐标 $(s,t)$，一个\trIncomingPoint{}$s=-\infty$ 与一个%
\trOutgoingPoint{}$s=+\infty$。
\end{itemize}
\trPointedBoundary\trRiemann\trSurface{}之间的\trIsomorphism{}是\trBiholomorphic{}映射，它们保持无穷远处的%
\trIncomingPoint{}与\trOutgoingPoint{}之间的区分。
$S$ 的一组\textbf{\trLagrange\trLabel}（Lagrangian labels）是一个族 $\{L_C\}$，由\trExact{}的%
\trLagrange\trSubmanifold{}$L_C\subset M$ 组成，以\trConnectedComponent{}$C\subset\partial S$ 标号。
每个 $\zeta\in\Sigma\subset\widehat{S}$ 都落在两个边界\trConnectedComponent{}$C_{\zeta,0}$ 与
$C_{\zeta,1}$（二者可以重合）的闭包之中；于是，在有一组\trLagrange\trLabel{}时，每个 $\zeta$ 都有一对
与之相配的\trLagrange\trSubmanifold{}$(L_{\zeta,0},L_{\zeta,1})$。我们的约定是：对 $\zeta\in\Sigma^-$，
沿自然边界\trOrientation，$C_{\zeta,1}$ 排在 $\zeta$ 之前，而 $C_{\zeta,0}$ 排在其后；对
$\zeta\in\Sigma^+$ 则恰好相反。
